%% kapa: exam=1 round=35 year=2024 period=1 form=A subject=civil_law subject_ko=민법 seq=31 number=31
%% source: Q-Net 감정평가사 기출문제 (cst00309) artlSeq=5213229 file=raw/1008/5213229/1교시 문제지 원본.pdf
%% conversion: kordoc 4.24.51 -> tools/gen1.py (auto); review=proofread-sonnet; flags=none
\begin{question}{31}
X토지를 3분의 1씩 공유하는 甲, 乙, 丙의 법률관계에 관한 설명으로 옳은 것은? (다툼이 있으면 판례에 따름)
\begin{choices}
  \item 甲이 乙과 丙의 동의 없이 X토지 중 3분의 1을 배타적으로 사용하는 경우, 乙은 방해배제를 청구할 수 없다.
  \item 甲과 乙이 협의하여 X토지를 매도하면 그 효력은 丙의 지분에도 미친다.
  \item 丁이 X토지의 점유를 무단으로 침해하고 있는 경우, 甲은 X토지 중 자신의 지분에 한하여 반환을 청구할 수 있다.
  \item 甲이 자신의 지분을 포기하더라도 乙과 丙이 이전등기를 하여야 甲의 지분을 취득한다.
  \item 丙이 1년 이상 X토지의 관리비용을 부담하지 않은 경우, 甲과 乙은 丙의 지분을 무상으로 취득할 수 있다.
\end{choices}
\end{question}
\answer{4}
